%!TEX root = ../../exa-ma-deliverables.tex
% D7.1 V3 / v2.90.0 -- WP7 CONTRIBUTION
%
% The content inherited from v2 is provided as a starting point.
% The WP leader/editor must:
% 1. review, validate, update or remove all inherited text;
% 2. identify explicitly what is new in 2026;
% 3. distinguish 2026 achievements from earlier baseline results;
% 4. connect the advances to Exa-MA software and applications;
% 5. associate applications with verification, validation or performance benchmarks;
% 6. report scientific, numerical and performance results;
% 7. provide software versions, machine configurations, references and evidence links;
% 8. identify missing application or benchmark mappings instead of inventing connections.

\chapter{WP7 -- Coordination and Activity Report (D7.4)}
\label{chap:wp7}

\section{Main advances in 2026}
\label{sec:wp7-main-advances-2026}
\WPUpdateNote{Complete one record for each main 2026 coordination, software,
application, benchmarking or training advance.  Keep the identifiers stable so
that contributions can be compared across annual reports.}

% Duplicate this record in WP7.tex for each main 2026 advance and replace the
% compact field cues with factual entries.  Do not edit the central table style.
\begin{WPAdvanceRecord}{WP7-A2026-1}
Exa-MA task and bottleneck & \WPFieldHint{T7.x; Bx} \\
Concise description of the 2026 advance & \WPFieldHint{What changed in 2026} \\
Responsible contributors and institutions & \WPFieldHint{Name; institution} \\
Implementing software & \WPFieldHint{Software; repository} \\
Software release, commit or PR & \WPFieldHint{Versioned identifier} \\
Associated application ID & \WPFieldHint{app-...} \\
Associated benchmark ID & \WPFieldHint{BM-... or no mapping yet} \\
Scientific or numerical result & \WPFieldHint{Value; unit; baseline} \\
Performance or scalability result & \WPFieldHint{Machine; resources; comparison} \\
Evidence and raw-data location & \WPFieldHint{URL; DOI; repository path} \\
Maturity & \WPFieldHint{method / prototype / integrated / benchmarked / reproducible} \\
Limitations and next step & \WPFieldHint{Limitation; next step} \\
Cross-WP contribution or dependency & \WPFieldHint{WP; role or dependency} \\
\end{WPAdvanceRecord}

\section{Co-design with PC5 and application demonstrators: initial feedback towards D7.2}
\label{sec:wp7:codesign-initial-feedback}

This section records an initial contribution to the T7.2 co-design report,
D7.2, scheduled at M48 and M60. It connects application needs to the
methodological and software contributions described elsewhere in this report.
The 2025 activity-report material records Exa-MA interactions with the PC5
working groups on high-order discretization, block-structured AMR and
unstructured meshing. The accounts below retain the scope of the documented
work and identify the feedback still needed for D7.2.
% Sources: Exa-MA ANR 2026 reporting material for 2025, section
% "Articulation du projet avec les autres projets de ce PEPR";
% WP1/T1.4; PC5 Ktirio Urban Building context; Feel++ implementation fiche.

\subsection{PC5 working group on efficient high-order discretization}

The WP1/T1.4 contribution describes a wave-equation proxy application
and a portable high-order finite-element library developed within the PC5
working group on Efficient High-Order Discretization. The proxy provides
an application setting for CPU/GPU kernels through Kokkos, with high-order
continuous and discontinuous Galerkin formulations. Its methodological scope
and development steps remain in WP1; the related FUnTiDES wave-propagation
software is described separately in \Cref{sec:FUnTiDES:software}.

The co-design link is the use of a wave application to guide numerical-kernel
and portability requirements. The WP1 account identifies face-term fusion,
additional finite-element spaces and transfer operators as further development
steps. This initial account does not add a performance comparison or infer
that these steps have been completed.

\subsection{PC5 working groups on block-structured AMR and unstructured meshing}

The recorded interactions with these two groups connect Exa-MA discretization
work to the PC5 application context. The existing technical material provides
separate descriptions of Cartesian adaptive-mesh methods in the Samurai fiche
(\Cref{sec:Samurai:software}), and of unstructured mesh adaptation and geometry
processing in WP1. These descriptions establish the available methods and
software; they do not, by themselves, establish decisions reached within the
working groups.

\todo[WP7 / D7.2 / PC5 AMR and unstructured-mesh feedback]{Document the
application requirements discussed with each working group, the resulting
methodological or software choices, and the remaining interface or validation
questions, with links to the corresponding discussion records or contributions.}

\subsection{Ktirio Urban Building and the HiDALGO2 Centre of Excellence}

The joint work on Ktirio Urban Building supplies a concrete co-design context
for urban building energy simulation and decision support. The demonstrator
needs terrain and building geometry, urban-microclimate/building-energy
coupling, and processing of large ensembles of building outputs. The component
roles and their relationships are described in
\Cref{sec:pc5:ktirio-urban-building}; this section records their implications
for the collaboration rather than repeating the fiches.

The geometry workflow connects application inputs to Ktirio-Geom and its
benchmark application, \texttt{app-feelpp-ktirio-geom}
(\Cref{sec:app:specs:app-feelpp-ktirio-geom}). Distributed geometric searches
motivate the Feel++ BVH contribution
(\Cref{sec:Feelpp:implementation-2026:bvh}). Coupling urban microclimate and
building energy models motivates interface field transfer through preCICE
and the Feel++--OpenFOAM participants
(\Cref{sec:Feelpp:implementation-2026:precice}). Memory and I/O constraints
in ensemble simulations motivate streaming statistical processing in Feel++,
on which the Ktirio Ensemble/UQ framework builds
(\Cref{sec:Feelpp:implementation-2026:stats}).

These connections identify how demonstrator requirements guide component
work across WP1, WP6 and WP7. The reported terrain timings characterize
Ktirio-Geom's terrain-generation workflow; they do not quantify the complete
coupled energy or ensemble/UQ workflow. The existing fiches retain the
component limitations and assessment needs. Further co-design feedback should
record interface choices and assess geometry, coupling and statistical
processing against their respective application requirements.

\section{Inherited WP7 material to update}
\WPUpdateNote{Review the inherited methodology, software and application chapters
in place.  Update the benchmark and training chapters only from versioned or
dated evidence; do not duplicate their detailed records here.}

The following WP7 chapters contain the retained v2 methodology, software and
application material, the reconstructed benchmark registry, and the integrated
training baseline.  Detailed application specifications remain in
\Cref{chap:applications}; WP chapters should summarize results and link to them.
